\begin{lemma}[Computed stored colors depend only on bounded history]
Let \(\mathcal D\) be \noderef{TashkinovSeriesData} with \(n\) rounds, let
\(m\leq n\), and put
\(\mathcal D^{\leq m}=\noderef{TashkinovSeriesDataPrefix}(m,\mathcal D)\).
For every \(a\in\operatorname{Fin}(m+1)\), if \(\iota_D(a)\) is the
canonical same-value inclusion in \(\operatorname{Fin}(n+1)\), then
\[
 \noderef{TashkinovStoredColors}(\mathcal D^{\leq m},a)
   =\noderef{TashkinovStoredColors}(\mathcal D,\iota_D(a)).
\]
\end{lemma}
\begin{proof}
If \(a=0\), both computed sets are empty by the defining zero-index branch of
\noderef{TashkinovStoredColors}.

Suppose \(a>0\).  Before subtracting missing colors, the left-hand union runs
over precisely those \(h\in\operatorname{Fin}(m+1)\) with \(h\leq a\).
Sending such an \(h\) through the canonical inclusion gives an old index of
the same numerical value, still at most \(\iota_D(a)\).  Conversely, if an old
index \(j\in\operatorname{Fin}(n+1)\) satisfies \(j\leq\iota_D(a)\), then
\(j\leq a<m+1\) numerically, so the same numerical value defines a unique
\(h\in\operatorname{Fin}(m+1)\).  These maps are inverse and the raw stored
sets agree at paired indices by the definition of
\noderef{TashkinovSeriesDataPrefix}.  Hence the two bounded unions are equal.

At the current index, the prefix coloring and tree are respectively the old
coloring at \(\iota_D(a)\) and the old tree at the same-value tree inclusion.
The latter is exactly the cast-successor of \(\iota_D(a)\).  Therefore the two
tree-missing-color sets subtracted from the unions are equal.  Equal unions
minus equal missing-color sets give the asserted equality.
\end{proof}
